\writeSectionFrame{\presentationTitle}{UML-Aktivitätsdiagramm}

\begin{frame}{Das Aktivitätsdiagramm}
	\begin{block}{beantwortet die Frage}
		Wie realisiert mein System ein bestimmtes Verhalten?
	\end{block}
\end{frame}

\begin{frame}{Beispiele für Aktivitäten}
	\begin{itemize}
	  \item Onlinebestellung
	  \item Waschvorgang von 90-Grad-Kochwäsche
	  \item Buchung einer bestandenen Prüfung
	  \item Bezahlvorgang
	  \item Kaffee kochen
	\end{itemize}
\end{frame}
\note{
	Eine Aktivität könnte z.B. eine Onlinebestellung vom ersten Klick bis zur Abgabe des Pakets an der Haustür sein. Ein Waschvorgang mit Berücksichtigung des Härtegrades des Wassers und der Umweltschutzbestimmungen ist eine Aktivität. Ebenso die Buchung einer bestandenen	Prüfung in einem Prüfungssystem mit Überprüfung aller Randbedingungen. Bezahlvorgänge 
	mit ebenfalls entsprechenden Prüfungen oder auch etwas so alltägliches wie Kaffee kochen. Auch das Aktivitätsdiagramm ist nicht auf das Modellieren technischer Probleme beschränkt. Alle beschriebenen Aktivitäten sind auf jeden Fall alles andere als trivial, wenn man alles miteinbedenkt.
}

\begin{frame}{Aktionen}
	\begin{center}
	  \includegraphics[scale=0.2]{\BILDERPFAD/aktivitaet1.png}
	\end{center}
\end{frame}
\note{
	Eine Aktion steht für den Aufruf eines Verhaltens oder die Bearbeitung von Daten. Dabei ist sie atomar, d.h. sie wird innerhalb dieser Aktivität nicht weiter zerlegt, könnte aber innerhalb einer	anderen Aktivität noch einmal genauer beschrieben werden. Die Aktivität ist ein Einzelschritt, der 
	zur Realisierung des durch die Aktivität beschriebenen Verhaltens beiträgt. Eine Aktivität ist das zentrale Element eines Aktivitätsdiagramms und alle anderen Elemente existieren nur, um die	Aktivitäten miteinander zu verknüpfen. Eine Aktion darf einen beliebigen Text enthalten. Dies kann
	ein Name sein, aber auch eine ausführlichere Beschreibung, wobei man da bei komplexen Abläufen vorsichtig sein sollte. Auch hier ist die Übersichtlichkeit am wichtigsten. 
}

\begin{frame}{Kanten}
	\begin{center}
	  \includegraphics[scale=0.2]{\BILDERPFAD/aktivitaet2.png}
	\end{center}
\end{frame}
\note{
	Kanten sind Übergänge zwischen zwei Knoten (Aktionen, Objektknoten usw). Kanten sind immer gerichtet und können mit einem Namen versehen werden. Die Kanten sind die Transportwege der Token. Damit eine Aktion ihre Arbeit aufnimmt, müssen an allen eingehenden Kanten gleichzeitig Token anliegen. Sobald eine Aktivität startet, werden die Token von den Kanten entfernt. Ist die Aktion abgeschlossen, bietet sie an allen ausgehenden Kanten jeweils ein Token an. Diese Token enthalten dann eine Eingabedaten für die nächste Aktion. 
}

\begin{frame}{Kontrollelemente}
	\begin{center}
	   \includegraphics[scale=0.25]{\BILDERPFAD/aktivitaet3.png}
	\end{center}
\end{frame}

\begin{frame}{Kontrollelemente}
	\begin{center}
	  \includegraphics[scale=0.3]{\BILDERPFAD/aktivitaet4.png}
	\end{center}
\end{frame}

\begin{frame}{Beispiel}
	\begin{center}
	  \includegraphics[scale=0.25]{\BILDERPFAD/aktivitaet5.png}
	\end{center}
\end{frame}

\begin{frame}{Startknoten}
	\begin{center}
	  \includegraphics[scale=0.4]{\BILDERPFAD/aktivitaet6.png}
	\end{center}
\end{frame}

\note{
	Ein Startknoten markiert den Startpunkt eines Ablaufs. Beim Start einer Aktivität wird ein Token generiert und an allen Kanten angeboten.
}

\begin{frame}{Endknoten}
	\begin{center}
	  \includegraphics[scale=0.4]{\BILDERPFAD/aktivitaet7.png}
	\end{center}
\end{frame}

\note{
    Der Endknoten stellt das Ende der Aktivität dar.
}

\begin{frame}{Verzweigungsknoten}
	\begin{center}
	  \includegraphics[scale=0.4]{\BILDERPFAD/aktivitaet8.png}
	\end{center}
\end{frame}

\note{
	   An einem Verzweigungsknoten wird eine Kante in mehrere Alternativen aufgespaltet, deren Ausführung jeweils von einer Bedingung abhängig ist. Ein Token passiert nur eine einzige ausgehende Kante. Bedingungen, sog. Guards, geben die Bedingungen an den ausgehenden Kanten an. Token müssen eindeutig geleitet werden können.
}

\begin{frame}{Beispiel}
	\begin{center}
	  \includegraphics[scale=0.3]{\BILDERPFAD/aktivitaet9.png}
	\end{center}
\end{frame}

\begin{frame}{Parallelierungsknoten}
	\begin{center}
	  \includegraphics[scale=0.5]{\BILDERPFAD/aktivitaet10.png}
	\end{center}
\end{frame}

\note{
	 An einem Parallelierungsknoten wird der eingehende Ablauf in mehrere parallele Abläufe aufgeteilt. 
}

\begin{frame}{Synchronisationsknoten}
	\begin{center}
	  \includegraphics[scale=0.3]{\BILDERPFAD/aktivitaet11.png}
	\end{center}
\end{frame}

\note{
	 Der Synchronisationsknoten führt parallele Abläufe wieder zusammen. Dazu müssen an allen Eingangskanten gleichzeitig Token angeboten werden. 
}

\begin{frame}{Aktivitätsbereiche (swimlanes)}
	\begin{center}
	  \includegraphics[scale=0.45]{\BILDERPFAD/aktivitaet12.png}
	\end{center}
\end{frame}

\note{
	Aktivitätsbereiche fassen Aktionen zusammen, die bestimmte Gemeinsamkeiten besitzen, z. B. organisatorische Einheiten oder Zuständigkeitsbereiche. Aktivitätsbereiche dürfen sowohl vertikal als auch horizontal modelliert werden und sich sogar schneiden.
}

\begin{frame}{Aktivitätsbereiche (swimlanes)}
	\begin{center}
	  \includegraphics[scale=0.35]{\BILDERPFAD/aktivitaet13.png}
	\end{center}
\end{frame}

\begin{frame}{Gast holt sich Essen oder Getränke}
	\begin{center}
	  \includegraphics[scale=0.23]{\BILDERPFAD/aktivitaet14.png}
	\end{center}
\end{frame}